\subsection{PID 制御}

リレー制御は偏差の符号だけで入力を切り替えるため実装しやすいが、目標付近のチャタリングや持続振動を生じやすい。比例制御は偏差の大きさに応じて入力を連続的に変えるが、一定外乱や摩擦に対する定常偏差が残る場合がある。PD 制御は変化率によって振動を抑える一方で測定ノイズに敏感であり、PI 制御は累積偏差で残留偏差を減らす一方で入力飽和時に積分状態が過大になりやすい。PID 制御は、これらの効果を一つの線形制御器としてまとめ、応答速度、定常偏差、振動抑制、実装上の制約を同時に調整するために導入する。

\begin{definition}[理想 PID 制御器]
制御器へ入力される偏差を $e(t)$ とする。測定値を使う単純なフィードバックでは $e(t)=r(t)-y_m(t)$ であり、測定ノイズを無視した設計式では $e(t)=r(t)-y(t)$ と書く。理想 PID 制御器は
\begin{equation}
  u_c(t)=K_P e(t)+K_I\int_0^t e(\tau)\,\dd\tau+K_D\dot{e}(t)
  \label{eq:pid-time}
\end{equation}
で計算入力 $u_c$ を作る制御器である。
\end{definition}

ラプラス変換を使うと、初期値を 0 として
\begin{align}
  U_c(s)&=K_P E(s)+K_I\frac{1}{s}E(s)+K_DsE(s),\\
  \frac{U_c(s)}{E(s)}&=K_P+\frac{K_I}{s}+K_Ds.
\end{align}
したがって PID 制御器の伝達関数は
\begin{equation}
  C(s)=K_P+\frac{K_I}{s}+K_Ds
\end{equation}
である。比例項は現在の偏差に作用し、積分項は偏差の累積に作用し、微分項は偏差の変化率に作用する。

実装では理想微分 $s$ をそのまま使うと高周波ノイズを増幅するため、
\begin{equation}
  K_D\frac{Ns}{s+N}
\end{equation}
のようなフィルタ付き微分を使う。目標値ステップに対する微分キックを避ける場合は、偏差全体ではなく測定値側だけに微分項を作用させる。入力に飽和がある場合、積分器が偏差を蓄積し続けるワインドアップが起こる。飽和前の計算入力を $u_c$、実入力を $u$ とし、積分状態を $z$ とすると、バック計算によるアンチワインドアップの一つは
\begin{equation}
  \dot{z}=e+k_{aw}(u-u_c)
\end{equation}
である。$u=u_c$ なら通常の積分 $\dot{z}=e$ に戻り、飽和中は差 $u-u_c$ によって積分状態を戻す。
